```latex
\documentclass[11pt]{article}

% ------------------------------------------------------------
% Packages
% ------------------------------------------------------------
\usepackage[margin=1in]{geometry}
\usepackage{amsmath}
\usepackage{amssymb}
\usepackage{booktabs}
\usepackage{array}
\usepackage{enumitem}
\usepackage{xcolor}
\usepackage{hyperref}

\hypersetup{
    colorlinks=true,
    linkcolor=blue,
    urlcolor=blue
}

% ------------------------------------------------------------
% Custom commands
% ------------------------------------------------------------
\newcommand{\vect}[1]{\mathbf{#1}}
\newcommand{\uvect}[1]{\hat{\mathbf{#1}}}

% ------------------------------------------------------------
% Document
% ------------------------------------------------------------
\begin{document}

% ============================================================
% TITLE
% ============================================================

\begin{titlepage}
    \centering

    \vspace*{2cm}

    {\Huge\bfseries Electromagnetics Foundations\par}

    \vspace{0.5cm}

    {\Large A Beginner-Friendly Study Guide\par}

    \vspace{1.5cm}

    {\large
    Vector Calculus, Electric Fields, Maxwell's Equations,\\
    Electromagnetic Waves, Fourier Analysis, and SDR Connections
    \par}

    \vfill

    {\large
    \textbf{Purpose:}\\
    Build the mathematical and physical foundation needed\\
    to understand graduate-level electromagnetics.
    \par}

    \vspace{1cm}

    {\large
    \textbf{Boxed equations = equations worth memorizing}
    \par}

    \vfill

    {\large
    \today
    \par}

\end{titlepage}

\tableofcontents
\newpage


% ============================================================
% 1. BIG PICTURE
% ============================================================

\section{The Big Picture}

Electromagnetics is fundamentally about understanding how
\textbf{electric fields} and \textbf{magnetic fields} are produced,
how they interact with matter, and how they propagate through space.

A useful mental hierarchy is:

\begin{enumerate}
    \item Scalars and vectors
    \item Vector calculus
    \item Electric charge and current
    \item Electric and magnetic fields
    \item Constitutive relations
    \item Maxwell's equations
    \item Electromagnetic waves
    \item Fourier analysis and frequency-domain thinking
\end{enumerate}

The most important conceptual chain is:

\[
\boxed{
\text{Charge/current}
\rightarrow
\text{fields}
\rightarrow
\text{Maxwell's equations}
\rightarrow
\text{waves}
\rightarrow
\text{energy/information propagation}
}
\]

For SDR work, Fourier analysis provides another important chain:

\[
\boxed{
\text{Time-domain signal}
\rightarrow
\text{Fourier transform}
\rightarrow
\text{frequency-domain representation}
\rightarrow
\text{filtering/modulation/demodulation}
}
\]


% ============================================================
% 2. SCALARS AND VECTORS
% ============================================================

\section{Scalars and Vectors}

A \textbf{scalar} has only magnitude.

Examples:

\begin{itemize}
    \item Temperature
    \item Voltage
    \item Electric charge
    \item Frequency
    \item Permittivity
\end{itemize}

A \textbf{vector} has both magnitude and direction.

Examples:

\begin{itemize}
    \item Electric field $\vect{E}$
    \item Magnetic field $\vect{H}$
    \item Current density $\vect{J}$
    \item Force $\vect{F}$
    \item Position $\vect{r}$
\end{itemize}

A three-dimensional vector can be written as

\[
\boxed{
\vect{A}
=
A_x\uvect{x}
+
A_y\uvect{y}
+
A_z\uvect{z}
}
\]

where:

\begin{itemize}
    \item $\vect{A}$ = vector
    \item $A_x,A_y,A_z$ = components of the vector
    \item $\uvect{x},\uvect{y},\uvect{z}$ = unit vectors in the $x$, $y$, and $z$ directions
\end{itemize}


% ============================================================
% 3. UNIT VECTORS
% ============================================================

\section{Unit Vectors}

A unit vector has magnitude 1.

The unit vector pointing in the direction of $\vect{A}$ is

\[
\boxed{
\hat{\vect{A}}
=
\frac{\vect{A}}{|\vect{A}|}
}
\]

The hat indicates that the vector has been normalized.

For example,

\[
\uvect{x}
\]

means ``the unit vector pointing in the $x$ direction.''


% ============================================================
% 4. VECTOR MAGNITUDE
% ============================================================

\section{Vector Magnitude}

The magnitude is the length of a vector.

For

\[
\vect{A}
=
A_x\uvect{x}
+
A_y\uvect{y}
+
A_z\uvect{z},
\]

the magnitude is

\[
\boxed{
|\vect{A}|
=
\sqrt{
A_x^2+A_y^2+A_z^2
}
}
\]

This is just the three-dimensional version of the Pythagorean theorem.


% ============================================================
% 5. DOT PRODUCT
% ============================================================

\section{Dot Product}

The dot product takes two vectors and produces a scalar.

\[
\boxed{
\vect{A}\cdot\vect{B}
=
|\vect{A}||\vect{B}|\cos\theta
}
\]

where $\theta$ is the angle between the vectors.

The dot product tells you how much one vector points in the direction of another.

Important cases:

\[
\vect{A}\cdot\vect{B}>0
\]

means the vectors generally point in similar directions.

\[
\vect{A}\cdot\vect{B}=0
\]

means the vectors are perpendicular.

\[
\vect{A}\cdot\vect{B}<0
\]

means they generally point in opposite directions.

The component form is

\[
\vect{A}\cdot\vect{B}
=
A_xB_x+A_yB_y+A_zB_z.
\]


% ============================================================
% 6. CROSS PRODUCT
% ============================================================

\section{Cross Product}

The cross product takes two vectors and produces another vector.

\[
\boxed{
\vect{A}\times\vect{B}
=
|\vect{A}||\vect{B}|
\sin\theta\,
\hat{\vect{n}}
}
\]

where:

\begin{itemize}
    \item $\theta$ = angle between the vectors
    \item $\hat{\vect{n}}$ = unit vector perpendicular to both
\end{itemize}

The cross product is zero when the vectors are parallel.

It is largest when the vectors are perpendicular.

Direction is determined using the right-hand rule.


% ============================================================
% 7. SCALAR AND VECTOR FIELDS
% ============================================================

\section{Scalar Fields and Vector Fields}

A \textbf{scalar field} assigns a scalar value to every point in space.

Example:

\[
V(x,y,z)
\]

is a voltage field.

A \textbf{vector field} assigns a vector to every point in space.

Example:

\[
\vect{E}(x,y,z)
\]

is an electric-field vector field.

This distinction is extremely important because vector calculus operators behave differently depending on whether their input is a scalar or vector.


% ============================================================
% 8. PARTIAL DERIVATIVES
% ============================================================

\section{Partial Derivatives}

A partial derivative tells you how a function changes with respect to one variable while treating the others as constant.

For

\[
V(x,y,z),
\]

we can calculate

\[
\frac{\partial V}{\partial x},
\qquad
\frac{\partial V}{\partial y},
\qquad
\frac{\partial V}{\partial z}.
\]

For example,

\[
V(x,y)=x^2+3y,
\]

gives

\[
\frac{\partial V}{\partial x}=2x
\]

and

\[
\frac{\partial V}{\partial y}=3.
\]


% ============================================================
% 9. NABLA
% ============================================================

\section{The Nabla Operator}

The symbol

\[
\nabla
\]

is called \textbf{nabla} or \textbf{del}.

In Cartesian coordinates,

\[
\boxed{
\nabla
=
\uvect{x}\frac{\partial}{\partial x}
+
\uvect{y}\frac{\partial}{\partial y}
+
\uvect{z}\frac{\partial}{\partial z}
}
\]

Think of $\nabla$ as a mathematical operator that performs spatial derivatives.

It can be combined with fields in several important ways:

\[
\nabla V
\]

gradient,

\[
\nabla\cdot\vect{A}
\]

divergence,

and

\[
\nabla\times\vect{A}
\]

curl.


% ============================================================
% 10. GRADIENT
% ============================================================

\section{Gradient}

The gradient takes a scalar field and produces a vector.

\[
\boxed{
\nabla V
=
\frac{\partial V}{\partial x}\uvect{x}
+
\frac{\partial V}{\partial y}\uvect{y}
+
\frac{\partial V}{\partial z}\uvect{z}
}
\]

The gradient points in the direction where the scalar field increases most rapidly.

For voltage,

\[
\boxed{
\vect{E}=-\nabla V
}
\]

This means the electric field points toward decreasing electric potential.


% ============================================================
% 11. DIVERGENCE
% ============================================================

\section{Divergence}

Divergence measures how much a vector field is ``spreading out'' from a point.

For

\[
\vect{A}
=
A_x\uvect{x}
+
A_y\uvect{y}
+
A_z\uvect{z},
\]

the divergence is

\[
\boxed{
\nabla\cdot\vect{A}
=
\frac{\partial A_x}{\partial x}
+
\frac{\partial A_y}{\partial y}
+
\frac{\partial A_z}{\partial z}
}
\]

The result is a scalar.

Intuitively:

\begin{itemize}
    \item Positive divergence $\rightarrow$ source
    \item Negative divergence $\rightarrow$ sink
    \item Zero divergence $\rightarrow$ no net source/sink
\end{itemize}


% ============================================================
% 12. CURL
% ============================================================

\section{Curl}

Curl measures local rotation or circulation of a vector field.

\[
\boxed{
\nabla\times\vect{A}
}
\]

For

\[
\vect{A}
=
A_x\uvect{x}
+
A_y\uvect{y}
+
A_z\uvect{z},
\]

the expanded expression is

\[
\nabla\times\vect{A}
=
\begin{vmatrix}
\uvect{x} & \uvect{y} & \uvect{z}\\
\frac{\partial}{\partial x}
&
\frac{\partial}{\partial y}
&
\frac{\partial}{\partial z}
\\
A_x&A_y&A_z
\end{vmatrix}.
\]

Therefore,

\[
\nabla\times\vect{A}
=
\left(
\frac{\partial A_z}{\partial y}
-
\frac{\partial A_y}{\partial z}
\right)\uvect{x}
+
\left(
\frac{\partial A_x}{\partial z}
-
\frac{\partial A_z}{\partial x}
\right)\uvect{y}
+
\left(
\frac{\partial A_y}{\partial x}
-
\frac{\partial A_x}{\partial y}
\right)\uvect{z}.
\]


% ============================================================
% 13. LAPLACIAN
% ============================================================

\section{Laplacian}

The Laplacian of a scalar field is

\[
\boxed{
\nabla^2V
=
\nabla\cdot(\nabla V)
}
\]

In Cartesian coordinates,

\[
\boxed{
\nabla^2V
=
\frac{\partial^2V}{\partial x^2}
+
\frac{\partial^2V}{\partial y^2}
+
\frac{\partial^2V}{\partial z^2}
}
\]

The Laplacian tells you how a quantity compares to its surrounding values.

It appears constantly in electromagnetics, especially in electrostatics and wave equations.


% ============================================================
% 14. INTEGRALS
% ============================================================

\section{Integrals}

An integral is essentially accumulation.

A one-dimensional integral,

\[
\int_a^b f(x)\,dx,
\]

adds up infinitely many small contributions.

In three-dimensional electromagnetics, you will encounter:

\[
\iiint_V
\]

for volume integrals,

\[
\iint_S
\]

for surface integrals,

and

\[
\oint_C
\]

for closed line integrals.

The subscript tells you what geometric object you are integrating over.


% ============================================================
% 15. DIFFERENTIAL DISPLACEMENT
% ============================================================

\section{Differential Displacement}

The differential displacement vector represents a tiny movement through space.

\[
\boxed{
d\vect{l}
=
dx\,\uvect{x}
+
dy\,\uvect{y}
+
dz\,\uvect{z}
}
\]

where:

\begin{itemize}
    \item $dx$ = tiny displacement in $x$
    \item $dy$ = tiny displacement in $y$
    \item $dz$ = tiny displacement in $z$
\end{itemize}


% ============================================================
% 16. DIFFERENTIAL SURFACE AREA
% ============================================================

\section{Differential Surface Area}

A small surface element is represented by

\[
\boxed{
d\vect{S}
=
\hat{\vect{n}}\,dS
}
\]

where:

\begin{itemize}
    \item $dS$ = tiny area
    \item $\hat{\vect{n}}$ = unit vector perpendicular to the surface
    \item $d\vect{S}$ = vector area element
\end{itemize}


% ============================================================
% 17. CLOSED INTEGRALS
% ============================================================

\section{Closed Integrals}

A circle over an integral sign means the path or surface is closed.

For example,

\[
\oint_C
\]

means a closed line integral around contour $C$.

Similarly,

\[
\oiint_S
\]

means a closed surface integral.

This notation becomes extremely important in Maxwell's equations.


% ============================================================
% 18. FLUX
% ============================================================

\section{Flux}

Flux describes how much of a vector field passes through a surface.

For a uniform field,

\[
\boxed{
\Phi
=
EA\cos\theta
}
\]

where:

\begin{itemize}
    \item $\Phi$ = flux
    \item $E$ = magnitude of the field
    \item $A$ = surface area
    \item $\theta$ = angle between the field and surface normal
\end{itemize}

The general electric-flux expression is

\[
\boxed{
\Phi_E
=
\int_S
\vect{E}\cdot d\vect{S}
}
\]

The dot product is important because only the component of $\vect{E}$ perpendicular to the surface contributes to flux.


% ============================================================
% 19. CHARGE
% ============================================================

\section{Charge}

Electric charge is measured in coulombs.

The total charge contained inside a volume is

\[
\boxed{
Q
=
\int_V\rho\,dv
}
\]

where:

\begin{itemize}
    \item $Q$ = total charge
    \item $\rho$ = volume charge density
    \item $dv$ = differential volume
\end{itemize}


% ============================================================
% 20. CHARGE DENSITY
% ============================================================

\section{Charge Density}

Charge can be distributed through a volume, across a surface, or along a line.

Volume charge density:

\[
\boxed{
\rho
=
\frac{dQ}{dv}
}
\]

Surface charge density:

\[
\boxed{
\rho_s
=
\frac{dQ}{dS}
}
\]

Line charge density:

\[
\boxed{
\rho_l
=
\frac{dQ}{dl}
}
\]

Units are:

\[
[\rho]=\mathrm{C/m^3}
\]

\[
[\rho_s]=\mathrm{C/m^2}
\]

\[
[\rho_l]=\mathrm{C/m}.
\]


% ============================================================
% 21. ELECTRIC POTENTIAL
% ============================================================

\section{Electric Potential}

Electric potential, or voltage, is measured in volts.

It represents electrical potential energy per unit charge.

The electric field and voltage are related by

\[
\boxed{
\vect{E}
=
-\nabla V
}
\]

The negative sign means that the electric field points toward decreasing voltage.

For a one-dimensional situation,

\[
E_x=-\frac{dV}{dx}.
\]


% ============================================================
% 22. ELECTRIC FIELD
% ============================================================

\section{Electric Field}

The electric field is force per unit charge.

\[
\boxed{
\vect{E}
=
\frac{\vect{F}}{q}
}
\]

where:

\begin{itemize}
    \item $\vect{E}$ = electric field
    \item $\vect{F}$ = electric force
    \item $q$ = charge
\end{itemize}

The units are

\[
\mathrm{N/C}
\]

or equivalently

\[
\mathrm{V/m}.
\]


% ============================================================
% 23. PERMITTIVITY
% ============================================================

\section{Permittivity}

Permittivity describes how a material responds to an electric field.

The symbol is

\[
\epsilon.
\]

The constitutive relation is

\[
\boxed{
\vect{D}
=
\epsilon\vect{E}
}
\]

where:

\begin{itemize}
    \item $\vect{D}$ = electric flux density
    \item $\epsilon$ = permittivity
    \item $\vect{E}$ = electric field
\end{itemize}

Permittivity is often written as

\[
\boxed{
\epsilon
=
\epsilon_r\epsilon_0
}
\]

where:

\begin{itemize}
    \item $\epsilon_r$ = relative permittivity
    \item $\epsilon_0$ = permittivity of free space
\end{itemize}

A material with larger $\epsilon$ generally stores more electric energy for a given electric field.


% ============================================================
% 24. DIELECTRICS
% ============================================================

\section{Dielectrics}

A dielectric is an electrically insulating material that can become polarized when exposed to an electric field.

Examples include:

\begin{itemize}
    \item Glass
    \item Ceramic
    \item Plastic
    \item Air
\end{itemize}

The relative permittivity is

\[
\epsilon_r
=
\frac{\epsilon}{\epsilon_0}.
\]

If

\[
\epsilon_r>1,
\]

the material has greater permittivity than free space.


% ============================================================
% 25. ELECTRIC FLUX DENSITY
% ============================================================

\section{Electric Flux Density}

Electric flux density is represented by

\[
\vect{D}.
\]

For a linear isotropic material,

\[
\boxed{
\vect{D}
=
\epsilon\vect{E}
}
\]

The units of $\vect{D}$ are

\[
\mathrm{C/m^2}.
\]

Do not confuse $\vect{D}$ with $\vect{E}$:

\[
\vect{E}
=
\text{electric field}
\]

while

\[
\vect{D}
=
\text{electric flux density}.
\]


% ============================================================
% 26. MAGNETIC FIELDS
% ============================================================

\section{Magnetic Flux Density and Magnetic Field Intensity}

There are two commonly used magnetic-field quantities:

\[
\vect{B}
\]

and

\[
\vect{H}.
\]

They are related through

\[
\boxed{
\vect{B}
=
\mu\vect{H}
}
\]

where:

\begin{itemize}
    \item $\vect{B}$ = magnetic flux density
    \item $\vect{H}$ = magnetic field intensity
    \item $\mu$ = permeability
\end{itemize}

Similarly,

\[
\boxed{
\mu
=
\mu_r\mu_0
}
\]

where:

\begin{itemize}
    \item $\mu_r$ = relative permeability
    \item $\mu_0$ = permeability of free space
\end{itemize}


% ============================================================
% 27. CURRENT DENSITY
% ============================================================

\section{Current Density}

Current density describes how electric current is distributed through an area.

It is represented by

\[
\vect{J}.
\]

For a simple ohmic material,

\[
\boxed{
\vect{J}
=
\sigma\vect{E}
}
\]

where:

\begin{itemize}
    \item $\vect{J}$ = current density
    \item $\sigma$ = electrical conductivity
    \item $\vect{E}$ = electric field
\end{itemize}

The units of $\vect{J}$ are

\[
\mathrm{A/m^2}.
\]


% ============================================================
% 28. CONSTITUTIVE RELATIONS
% ============================================================

\section{The Three Constitutive Relations}

For simple linear isotropic materials, remember:

\[
\boxed{
\vect{D}
=
\epsilon\vect{E}
}
\]

\[
\boxed{
\vect{B}
=
\mu\vect{H}
}
\]

\[
\boxed{
\vect{J}
=
\sigma\vect{E}
}
\]

These describe how a material responds to electric and magnetic fields.


% ============================================================
% 29. MAXWELL
% ============================================================

\section{Maxwell's Four Equations}

These are the central equations of classical electromagnetics.

\subsection{Gauss's Law for Electricity}

\[
\boxed{
\nabla\cdot\vect{D}
=
\rho
}
\]

Electric charge acts as a source or sink of electric flux.

\subsection{Gauss's Law for Magnetism}

\[
\boxed{
\nabla\cdot\vect{B}
=
0
}
\]

There are no isolated magnetic monopoles in classical electromagnetics.

\subsection{Faraday's Law}

\[
\boxed{
\nabla\times\vect{E}
=
-\frac{\partial\vect{B}}{\partial t}
}
\]

A changing magnetic field produces a circulating electric field.

\subsection{Ampere-Maxwell Law}

\[
\boxed{
\nabla\times\vect{H}
=
\vect{J}
+
\frac{\partial\vect{D}}{\partial t}
}
\]

Electric current and changing electric fields produce circulating magnetic fields.


% ============================================================
% 30. PHYSICAL PICTURE
% ============================================================

\section{Maxwell's Equations: Physical Picture}

A useful memorization table is:

\begin{center}
\begin{tabular}{lll}
\toprule
Equation & Mathematical operation & Physical idea \\
\midrule
$\nabla\cdot\vect{D}=\rho$
& Divergence
& Charge creates electric flux \\

$\nabla\cdot\vect{B}=0$
& Divergence
& No magnetic monopoles \\

$\nabla\times\vect{E}=-\partial\vect{B}/\partial t$
& Curl
& Changing magnetic field creates electric field \\

$\nabla\times\vect{H}=\vect{J}+\partial\vect{D}/\partial t$
& Curl
& Current/changing electric field creates magnetic field \\
\bottomrule
\end{tabular}
\end{center}

The deepest idea is that changing electric and magnetic fields can sustain each other and propagate through space.


% ============================================================
% 31. POISSON
% ============================================================

\section{Poisson's Equation}

Start with

\[
\vect{E}=-\nabla V.
\]

Using

\[
\vect{D}=\epsilon\vect{E},
\]

we get

\[
\vect{D}
=
-\epsilon\nabla V.
\]

Gauss's law says

\[
\nabla\cdot\vect{D}
=
\rho.
\]

Therefore,

\[
\nabla\cdot(-\epsilon\nabla V)
=
\rho.
\]

For constant $\epsilon$,

\[
-\epsilon\nabla^2V
=
\rho.
\]

Therefore,

\[
\boxed{
\nabla^2V
=
-\frac{\rho}{\epsilon}
}
\]

This is \textbf{Poisson's equation}.

It tells you how charge density determines the curvature of voltage.


% ============================================================
% 32. LAPLACE
% ============================================================

\section{Laplace's Equation}

If there is no charge,

\[
\rho=0.
\]

Poisson's equation becomes

\[
\boxed{
\nabla^2V=0
}
\]

This is \textbf{Laplace's equation}.

A huge amount of electrostatics involves solving Laplace's equation subject to boundary conditions.


% ============================================================
% 33. BOUNDARY CONDITIONS
% ============================================================

\section{Boundary Conditions}

When an electromagnetic field crosses a boundary between two materials, the field components must satisfy specific conditions.

For the normal component of $\vect{D}$,

\[
\boxed{
\hat{\vect{n}}
\cdot
(\vect{D}_2-\vect{D}_1)
=
\rho_s
}
\]

where:

\begin{itemize}
    \item $\hat{\vect{n}}$ = unit normal pointing from region 1 to region 2
    \item $\vect{D}_1$ = electric flux density on side 1
    \item $\vect{D}_2$ = electric flux density on side 2
    \item $\rho_s$ = surface charge density
\end{itemize}

For the tangential electric field,

\[
\boxed{
\hat{\vect{n}}
\times
(\vect{E}_2-\vect{E}_1)
=
0
}
\]

This says the tangential component of $\vect{E}$ is continuous across an ordinary static boundary.


% ============================================================
% 34. ELECTROMAGNETIC WAVES
% ============================================================

\section{Electromagnetic Waves}

A simple plane electromagnetic wave can be represented as

\[
\boxed{
\vect{E}
=
\uvect{x}
E_0
\cos(\omega t-kz)
}
\]

where:

\begin{itemize}
    \item $E_0$ = electric-field amplitude
    \item $\omega$ = angular frequency
    \item $t$ = time
    \item $k$ = wavenumber
    \item $z$ = position
\end{itemize}

The argument

\[
\omega t-kz
\]

determines the phase of the wave.

The electric field points in the $x$ direction while the wave propagates in the $+z$ direction.


% ============================================================
% 35. FREQUENCY
% ============================================================

\section{Frequency, Period, and Angular Frequency}

Frequency is measured in hertz.

\[
\boxed{
T=\frac{1}{f}
}
\]

where:

\begin{itemize}
    \item $T$ = period in seconds
    \item $f$ = frequency in hertz
\end{itemize}

Angular frequency is

\[
\boxed{
\omega=2\pi f
}
\]

where:

\begin{itemize}
    \item $\omega$ = angular frequency in radians/second
    \item $f$ = frequency in hertz
\end{itemize}


% ============================================================
% 36. WAVELENGTH
% ============================================================

\section{Wavelength and Wavenumber}

Wavelength is the physical distance corresponding to one complete cycle.

It is represented by

\[
\lambda.
\]

The wavenumber is

\[
\boxed{
k=\frac{2\pi}{\lambda}
}
\]

where:

\begin{itemize}
    \item $k$ = wavenumber in radians/meter
    \item $\lambda$ = wavelength in meters
\end{itemize}

The factor $2\pi$ appears because one complete cycle corresponds to $2\pi$ radians.


% ============================================================
% 37. WAVE VELOCITY
% ============================================================

\section{Wave Velocity}

Wave velocity is related to frequency and wavelength by

\[
\boxed{
v=f\lambda
}
\]

Since

\[
\omega=2\pi f
\]

and

\[
k=\frac{2\pi}{\lambda},
\]

we can also write

\[
\boxed{
v=\frac{\omega}{k}
}
\]

The key relationships are

\[
\boxed{
T=\frac{1}{f},
\qquad
\omega=2\pi f,
\qquad
k=\frac{2\pi}{\lambda},
\qquad
v=f\lambda,
\qquad
v=\frac{\omega}{k}
}
\]


% ============================================================
% 38. EM WAVE VELOCITY
% ============================================================

\section{Electromagnetic Wave Velocity}

In a homogeneous medium,

\[
\boxed{
v
=
\frac{1}{\sqrt{\mu\epsilon}}
}
\]

where:

\begin{itemize}
    \item $v$ = electromagnetic wave velocity
    \item $\mu$ = permeability
    \item $\epsilon$ = permittivity
\end{itemize}

In free space,

\[
\boxed{
c
=
\frac{1}{\sqrt{\mu_0\epsilon_0}}
}
\]

where $c$ is the speed of light.


% ============================================================
% 39. INTRINSIC IMPEDANCE
% ============================================================

\section{Intrinsic Impedance}

The ratio of electric-field magnitude to magnetic-field magnitude in a uniform plane wave is the intrinsic impedance.

\[
\boxed{
\eta
=
\sqrt{\frac{\mu}{\epsilon}}
}
\]

and

\[
\boxed{
\frac{E}{H}
=
\eta
}
\]

For free space,

\[
\boxed{
\eta_0
\approx
377~\Omega
}
\]

This is analogous to an impedance relationship between voltage and current, but here it describes the ratio between electric and magnetic field strengths.


% ============================================================
% 40. POYNTING VECTOR
% ============================================================

\section{Poynting Vector}

The Poynting vector describes electromagnetic power flow.

\[
\boxed{
\vect{S}
=
\vect{E}\times\vect{H}
}
\]

where:

\begin{itemize}
    \item $\vect{S}$ = power density
    \item $\vect{E}$ = electric field
    \item $\vect{H}$ = magnetic field
\end{itemize}

The units are

\[
\mathrm{W/m^2}.
\]

The direction of $\vect{S}$ tells you the direction in which electromagnetic energy is flowing.


% ============================================================
% 41. COMPLEX NUMBERS
% ============================================================

\section{Complex Numbers}

Electrical engineering frequently uses complex numbers.

The imaginary unit is

\[
\boxed{
j=\sqrt{-1}
}
\]

Electrical engineers use $j$ instead of $i$ because $i$ commonly represents current.

Euler's formula is

\[
\boxed{
e^{j\theta}
=
\cos\theta
+
j\sin\theta
}
\]

This is extremely important because it connects exponentials with sinusoidal signals.


% ============================================================
% 42. PHASORS
% ============================================================

\section{Phasors}

Suppose a sinusoidal signal is

\[
E(t)
=
E_0\cos(\omega t+\phi).
\]

We can represent its amplitude and phase using a complex phasor:

\[
\boxed{
\tilde{E}
=
E_0e^{j\phi}
}
\]

The physical signal is recovered by taking the real part:

\[
E(t)
=
\operatorname{Re}
\left\{
\tilde{E}e^{j\omega t}
\right\}.
\]

One of the major benefits of phasors is that differentiation with respect to time becomes multiplication by $j\omega$:

\[
\boxed{
\frac{\partial}{\partial t}
\longrightarrow
j\omega
}
\]

Therefore,

\[
\boxed{
\frac{\partial E}{\partial t}
\longrightarrow
j\omega\tilde{E}
}
\]

This makes many electromagnetic equations much easier to solve in the frequency domain.


% ============================================================
% 43. FOURIER TRANSFORM
% ============================================================

\section{Fourier Transform}

The Fourier transform converts a time-domain signal into its frequency-domain representation.

\[
\boxed{
X(f)
=
\int_{-\infty}^{\infty}
x(t)e^{-j2\pi ft}\,dt
}
\]

where:

\begin{itemize}
    \item $x(t)$ = time-domain signal
    \item $X(f)$ = frequency-domain representation
    \item $f$ = frequency
    \item $t$ = time
\end{itemize}

The inverse transform is

\[
x(t)
=
\int_{-\infty}^{\infty}
X(f)e^{j2\pi ft}\,df.
\]

The big idea is:

\[
\boxed{
\text{time domain}
\longleftrightarrow
\text{frequency domain}
}
\]


% ============================================================
% 44. DFT VS FFT
% ============================================================

\section{DFT vs FFT}

The Discrete Fourier Transform is

\[
\boxed{
X[k]
=
\sum_{n=0}^{N-1}
x[n]
e^{-j2\pi kn/N}
}
\]

where:

\begin{itemize}
    \item $x[n]$ = input samples
    \item $X[k]$ = frequency-bin output
    \item $N$ = number of samples
    \item $n$ = time-domain sample index
    \item $k$ = frequency-bin index
\end{itemize}

The DFT is the mathematical operation.

The FFT is an efficient algorithm for computing the DFT.

In simple terms:

\[
\boxed{
\text{FFT}
=
\text{fast algorithm for computing the DFT}
}
\]

For a naive DFT, computational complexity is approximately

\[
O(N^2).
\]

For a typical FFT algorithm,

\[
O(N\log N).
\]

This is why FFTs are so useful for SDR.


% ============================================================
% 45. SAMPLING
% ============================================================

\section{Sampling}

When converting an analog signal into digital samples, the sampling rate must be sufficiently high to represent the signal without aliasing.

The Nyquist condition is

\[
\boxed{
f_s>2f_{\max}
}
\]

where:

\begin{itemize}
    \item $f_s$ = sampling frequency
    \item $f_{\max}$ = highest frequency component of the signal
\end{itemize}

If the condition is violated, frequency components can alias into other frequencies.

In SDR, this is one of the fundamental reasons that sample rate and filtering matter.


% ============================================================
% 46. CONVOLUTION
% ============================================================

\section{Convolution}

Convolution describes how a linear time-invariant system transforms an input signal.

In the time domain,

\[
\boxed{
y(t)
=
x(t)*h(t)
}
\]

where:

\begin{itemize}
    \item $x(t)$ = input
    \item $h(t)$ = impulse response
    \item $y(t)$ = output
\end{itemize}

In the frequency domain, convolution becomes multiplication:

\[
\boxed{
Y(f)
=
X(f)H(f)
}
\]

This is one of the major reasons Fourier analysis is so useful for filters.

For SDR:

\[
\boxed{
\text{filtering in time}
\longleftrightarrow
\text{multiplication in frequency}
}
\]


% ============================================================
% 47. DIVERGENCE THEOREM
% ============================================================

\section{Divergence Theorem}

The divergence theorem connects a volume integral to a closed surface integral.

\[
\boxed{
\iiint_V
\nabla\cdot\vect{A}\,dv
=
\oiint_S
\vect{A}\cdot d\vect{S}
}
\]

In words:

\begin{quote}
The total amount of divergence inside a volume equals the net outward flux through its enclosing surface.
\end{quote}

This theorem is essential for understanding the integral form of Gauss's law.


% ============================================================
% 48. STOKES' THEOREM
% ============================================================

\section{Stokes' Theorem}

Stokes' theorem connects curl over a surface to circulation around the boundary.

\[
\boxed{
\iint_S
(\nabla\times\vect{A})
\cdot d\vect{S}
=
\oint_C
\vect{A}\cdot d\vect{l}
}
\]

In words:

\begin{quote}
The total curl passing through a surface equals the circulation around the boundary of that surface.
\end{quote}

This theorem is essential for understanding the integral form of Faraday's law and Ampere's law.


% ============================================================
% 49. DIFFERENTIAL EQUATIONS
% ============================================================

\section{Differential Equations}

A differential equation contains derivatives of an unknown function.

For example,

\[
\frac{d^2y}{dx^2}+ky=0.
\]

A common solution is sinusoidal:

\[
y(x)
=
A\cos(kx)
+
B\sin(kx).
\]

Electromagnetic fields are governed by differential equations because Maxwell's equations describe how fields vary across space and time.

The wave equation is a particularly important example.


% ============================================================
% 50. SEPARATION OF VARIABLES
% ============================================================

\section{Separation of Variables}

Separation of variables is a technique for solving certain partial differential equations.

For example, assume

\[
\boxed{
V(x,y)
=
X(x)Y(y)
}
\]

where:

\begin{itemize}
    \item $V(x,y)$ = unknown two-dimensional function
    \item $X(x)$ = function depending only on $x$
    \item $Y(y)$ = function depending only on $y$
\end{itemize}

The goal is to turn one complicated PDE into multiple simpler ordinary differential equations.


% ============================================================
% 51. MASTER EQUATION SHEET
% ============================================================

\section{Master Equation Sheet}

These are the equations to know especially well.

\subsection*{Vectors}

\[
\boxed{
|\vect{A}|
=
\sqrt{A_x^2+A_y^2+A_z^2}
}
\]

\[
\boxed{
\hat{\vect{A}}
=
\frac{\vect{A}}{|\vect{A}|}
}
\]

\[
\boxed{
\vect{A}\cdot\vect{B}
=
|\vect{A}||\vect{B}|\cos\theta
}
\]

\[
\boxed{
\vect{A}\times\vect{B}
=
|\vect{A}||\vect{B}|\sin\theta\,\hat{\vect{n}}
}
\]

\subsection*{Vector Calculus}

\[
\boxed{
\nabla V
=
\frac{\partial V}{\partial x}\uvect{x}
+
\frac{\partial V}{\partial y}\uvect{y}
+
\frac{\partial V}{\partial z}\uvect{z}
}
\]

\[
\boxed{
\nabla\cdot\vect{A}
=
\frac{\partial A_x}{\partial x}
+
\frac{\partial A_y}{\partial y}
+
\frac{\partial A_z}{\partial z}
}
\]

\[
\boxed{
\nabla\times\vect{A}
}
\]

\[
\boxed{
\nabla^2V
=
\nabla\cdot(\nabla V)
}
\]

\subsection*{Electrostatics}

\[
\boxed{
\vect{E}
=
-\nabla V
}
\]

\[
\boxed{
\vect{D}
=
\epsilon\vect{E}
}
\]

\[
\boxed{
Q
=
\int_V\rho\,dv
}
\]

\[
\boxed{
\nabla^2V
=
-\frac{\rho}{\epsilon}
}
\]

\[
\boxed{
\nabla^2V=0
\qquad
(\rho=0)
}
\]

\subsection*{Magnetics}

\[
\boxed{
\vect{B}
=
\mu\vect{H}
}
\]

\[
\boxed{
\vect{J}
=
\sigma\vect{E}
}
\]

\subsection*{Maxwell}

\[
\boxed{
\nabla\cdot\vect{D}
=
\rho
}
\]

\[
\boxed{
\nabla\cdot\vect{B}
=
0
}
\]

\[
\boxed{
\nabla\times\vect{E}
=
-\frac{\partial\vect{B}}{\partial t}
}
\]

\[
\boxed{
\nabla\times\vect{H}
=
\vect{J}
+
\frac{\partial\vect{D}}{\partial t}
}
\]

\subsection*{Waves}

\[
\boxed{
v
=
\frac{1}{\sqrt{\mu\epsilon}}
}
\]

\[
\boxed{
c
=
\frac{1}{\sqrt{\mu_0\epsilon_0}}
}
\]

\[
\boxed{
v=f\lambda
}
\]

\[
\boxed{
v=\frac{\omega}{k}
}
\]

\[
\boxed{
k=\frac{2\pi}{\lambda}
}
\]

\[
\boxed{
\omega=2\pi f
}
\]

\[
\boxed{
\eta
=
\sqrt{\frac{\mu}{\epsilon}}
}
\]

\[
\boxed{
\vect{S}
=
\vect{E}\times\vect{H}
}
\]

\subsection*{Complex Numbers and Phasors}

\[
\boxed{
j=\sqrt{-1}
}
\]

\[
\boxed{
e^{j\theta}
=
\cos\theta+j\sin\theta
}
\]

\[
\boxed{
\frac{\partial}{\partial t}
\longrightarrow
j\omega
}
\]

\subsection*{Fourier and SDR}

\[
\boxed{
X(f)
=
\int_{-\infty}^{\infty}
x(t)e^{-j2\pi ft}\,dt
}
\]

\[
\boxed{
X[k]
=
\sum_{n=0}^{N-1}
x[n]e^{-j2\pi kn/N}
}
\]

\[
\boxed{
f_s>2f_{\max}
}
\]

\[
\boxed{
Y(f)
=
X(f)H(f)
}
\]

\subsection*{Vector Integral Theorems}

\[
\boxed{
\iiint_V
\nabla\cdot\vect{A}\,dv
=
\oiint_S
\vect{A}\cdot d\vect{S}
}
\]

\[
\boxed{
\iint_S
(\nabla\times\vect{A})
\cdot d\vect{S}
=
\oint_C
\vect{A}\cdot d\vect{l}
}
\]


% ============================================================
% 52. SYMBOL REFERENCE
% ============================================================

\section{Symbol Reference}

\begin{center}
\begin{tabular}{lll}
\toprule
Symbol & Meaning & Typical Unit \\
\midrule
$\vect{E}$ & Electric field & V/m \\
$\vect{D}$ & Electric flux density & C/m$^2$ \\
$\vect{B}$ & Magnetic flux density & T \\
$\vect{H}$ & Magnetic field intensity & A/m \\
$\vect{J}$ & Current density & A/m$^2$ \\
$V$ & Electric potential & V \\
$Q$ & Electric charge & C \\
$\rho$ & Volume charge density & C/m$^3$ \\
$\rho_s$ & Surface charge density & C/m$^2$ \\
$\rho_l$ & Line charge density & C/m \\
$\epsilon$ & Permittivity & F/m \\
$\mu$ & Permeability & H/m \\
$\sigma$ & Conductivity & S/m \\
$\eta$ & Intrinsic impedance & $\Omega$ \\
$f$ & Frequency & Hz \\
$T$ & Period & s \\
$\omega$ & Angular frequency & rad/s \\
$\lambda$ & Wavelength & m \\
$k$ & Wavenumber & rad/m \\
$c$ & Speed of light & m/s \\
$\vect{S}$ & Poynting vector & W/m$^2$ \\
\bottomrule
\end{tabular}
\end{center}


% ============================================================
% 53. WHAT TO MEMORIZE FIRST
% ============================================================

\section{What to Memorize First}

Do not try to memorize everything simultaneously.

Start with these.

\subsection*{1. Vector Calculus}

Know what these mean:

\[
\boxed{
\nabla V
}
\]

\[
\boxed{
\nabla\cdot\vect{A}
}
\]

\[
\boxed{
\nabla\times\vect{A}
}
\]

\[
\boxed{
\nabla^2V
}
\]

Remember:

\[
\boxed{
\text{gradient}
\rightarrow
\text{scalar to vector}
}
\]

\[
\boxed{
\text{divergence}
\rightarrow
\text{vector to scalar}
}
\]

\[
\boxed{
\text{curl}
\rightarrow
\text{vector to vector}
}
\]

\subsection*{2. Electric Field and Voltage}

\[
\boxed{
\vect{E}
=
-\nabla V
}
\]

\subsection*{3. Constitutive Relations}

\[
\boxed{
\vect{D}=\epsilon\vect{E}
}
\]

\[
\boxed{
\vect{B}=\mu\vect{H}
}
\]

\[
\boxed{
\vect{J}=\sigma\vect{E}
}
\]

\subsection*{4. Maxwell's Equations}

Memorize all four.

\subsection*{5. Wave Relationships}

\[
\boxed{
v=f\lambda
}
\]

\[
\boxed{
\omega=2\pi f
}
\]

\[
\boxed{
k=\frac{2\pi}{\lambda}
}
\]

\[
\boxed{
v=\frac{\omega}{k}
}
\]

\subsection*{6. Fourier Analysis}

Understand the idea that

\[
\boxed{
\text{time domain}
\longleftrightarrow
\text{frequency domain}
}
\]

and know the DFT/FFT relationship.


% ============================================================
% 54. MOST IMPORTANT MENTAL MAP
% ============================================================

\section{The Most Important Mental Map}

Try to understand electromagnetics as one connected system.

\subsection*{Step 1: Charge}

Charge density creates electric flux:

\[
\boxed{
\nabla\cdot\vect{D}
=
\rho
}
\]

\subsection*{Step 2: Electric Field}

Electric field is related to voltage:

\[
\boxed{
\vect{E}
=
-\nabla V
}
\]

\subsection*{Step 3: Material}

The material determines how $\vect{D}$ relates to $\vect{E}$:

\[
\boxed{
\vect{D}
=
\epsilon\vect{E}
}
\]

\subsection*{Step 4: Current}

Conducting materials allow current:

\[
\boxed{
\vect{J}
=
\sigma\vect{E}
}
\]

\subsection*{Step 5: Magnetic Field}

Magnetic fields are related to currents and changing electric fields:

\[
\boxed{
\nabla\times\vect{H}
=
\vect{J}
+
\frac{\partial\vect{D}}{\partial t}
}
\]

\subsection*{Step 6: Changing Magnetic Fields}

Changing magnetic fields create electric fields:

\[
\boxed{
\nabla\times\vect{E}
=
-\frac{\partial\vect{B}}{\partial t}
}
\]

\subsection*{Step 7: Waves}

The electric and magnetic fields can therefore continually generate one another.

That produces electromagnetic waves.

\[
\boxed{
\text{changing }\vect{E}
\rightarrow
\text{changing }\vect{H}
\rightarrow
\text{changing }\vect{E}
\rightarrow
\cdots
}
\]


% ============================================================
% 55. MAXWELL AS A MENTAL MAP
% ============================================================

\section{Maxwell's Equations as a Mental Map}

Instead of memorizing Maxwell's equations as four unrelated equations, associate each with a question.

\subsection*{Question 1: What creates electric flux?}

\[
\boxed{
\nabla\cdot\vect{D}
=
\rho
}
\]

Answer:

\[
\text{electric charge}.
\]

\subsection*{Question 2: What creates magnetic flux divergence?}

\[
\boxed{
\nabla\cdot\vect{B}
=
0
}
\]

Answer:

\[
\text{there are no isolated magnetic monopoles}.
\]

\subsection*{Question 3: What creates circulating electric fields?}

\[
\boxed{
\nabla\times\vect{E}
=
-\frac{\partial\vect{B}}{\partial t}
}
\]

Answer:

\[
\text{changing magnetic fields}.
\]

\subsection*{Question 4: What creates circulating magnetic fields?}

\[
\boxed{
\nabla\times\vect{H}
=
\vect{J}
+
\frac{\partial\vect{D}}{\partial t}
}
\]

Answer:

\[
\text{electric current and changing electric fields}.
\]


% ============================================================
% 56. THREE QUESTIONS
% ============================================================

\section{Three Questions to Ask for Every Equation}

When looking at a new electromagnetics equation, ask:

\subsection*{1. What goes into the equation?}

For example,

\[
\nabla\cdot\vect{D}
\]

takes a vector field and produces a scalar.

\subsection*{2. What comes out?}

For example,

\[
\nabla\cdot\vect{D}=\rho
\]

produces charge density.

\subsection*{3. What does it physically mean?}

In this case:

\[
\boxed{
\text{electric charge is a source of electric flux}
}
\]

This method is often more useful than trying to memorize symbols without understanding them.


% ============================================================
% 57. FINAL STUDY STRATEGY
% ============================================================

\section{Final Study Strategy}

A useful order for learning the material is:

\begin{enumerate}
    \item Scalars vs vectors
    \item Unit vectors
    \item Vector magnitude
    \item Dot product
    \item Cross product
    \item Partial derivatives
    \item Gradient
    \item Divergence
    \item Curl
    \item Laplacian
    \item Line/surface/volume integrals
    \item Flux
    \item Charge density
    \item Electric potential
    \item Electric field
    \item Permittivity
    \item $\vect{D}$
    \item $\vect{B}$ and $\vect{H}$
    \item Current density
    \item Maxwell's equations
    \item Poisson and Laplace equations
    \item Boundary conditions
    \item Waves
    \item Phasors
    \item Fourier transforms
    \item DFT/FFT
    \item Sampling
    \item Convolution
\end{enumerate}

For each equation, learn three things:

\[
\boxed{
\text{What does it say?}
}
\]

\[
\boxed{
\text{What does every symbol mean?}
}
\]

\[
\boxed{
\text{What does it mean physically?}
}
\]

Do not worry about memorizing every derivation immediately.

The goal is first to understand the structure.


% ============================================================
% 58. QUICK REFERENCE
% ============================================================

\section{Quick Reference}

\subsection*{Vector Operations}

\[
\boxed{
|\vect{A}|
=
\sqrt{A_x^2+A_y^2+A_z^2}
}
\]

\[
\boxed{
\vect{A}\cdot\vect{B}
=
|\vect{A}||\vect{B}|\cos\theta
}
\]

\[
\boxed{
\vect{A}\times\vect{B}
=
|\vect{A}||\vect{B}|\sin\theta\,\hat{\vect{n}}
}
\]

\[
\boxed{
\nabla V
=
\text{gradient}
}
\]

\[
\boxed{
\nabla\cdot\vect{A}
=
\text{divergence}
}
\]

\[
\boxed{
\nabla\times\vect{A}
=
\text{curl}
}
\]

\[
\boxed{
\nabla^2V
=
\text{Laplacian}
}
\]

\subsection*{Electrostatics}

\[
\boxed{
\vect{E}
=
-\nabla V
}
\]

\[
\boxed{
\vect{D}
=
\epsilon\vect{E}
}
\]

\[
\boxed{
\nabla\cdot\vect{D}
=
\rho
}
\]

\[
\boxed{
\nabla^2V
=
-\frac{\rho}{\epsilon}
}
\]

\[
\boxed{
\nabla^2V=0
\quad
\text{if }
\rho=0
}
\]

\subsection*{Magnetics}

\[
\boxed{
\vect{B}
=
\mu\vect{H}
}
\]

\[
\boxed{
\nabla\cdot\vect{B}
=
0
}
\]

\subsection*{Maxwell}

\[
\boxed{
\nabla\cdot\vect{D}
=
\rho
}
\]

\[
\boxed{
\nabla\cdot\vect{B}
=
0
}
\]

\[
\boxed{
\nabla\times\vect{E}
=
-\frac{\partial\vect{B}}{\partial t}
}
\]

\[
\boxed{
\nabla\times\vect{H}
=
\vect{J}
+
\frac{\partial\vect{D}}{\partial t}
}
\]

\subsection*{Waves}

\[
\boxed{
v
=
\frac{1}{\sqrt{\mu\epsilon}}
}
\]

\[
\boxed{
v=f\lambda
}
\]

\[
\boxed{
v=\frac{\omega}{k}
}
\]

\[
\boxed{
\omega=2\pi f
}
\]

\[
\boxed{
k=\frac{2\pi}{\lambda}
}
\]

\[
\boxed{
\eta
=
\sqrt{\frac{\mu}{\epsilon}}
}
\]

\[
\boxed{
\vect{S}
=
\vect{E}\times\vect{H}
}
\]

\subsection*{Fourier / SDR}

\[
\boxed{
X(f)
=
\int_{-\infty}^{\infty}
x(t)e^{-j2\pi ft}\,dt
}
\]

\[
\boxed{
X[k]
=
\sum_{n=0}^{N-1}
x[n]e^{-j2\pi kn/N}
}
\]

\[
\boxed{
f_s>2f_{\max}
}
\]

\[
\boxed{
Y(f)
=
X(f)H(f)
}
\]


% ============================================================
% 59. CORE MENTAL MODEL
% ============================================================

\section{Core Mental Model}

If everything above feels overwhelming, remember this simplified model:

\[
\boxed{
\text{Charge}
\rightarrow
\text{Electric Field}
}
\]

\[
\boxed{
\text{Current}
\rightarrow
\text{Magnetic Field}
}
\]

\[
\boxed{
\text{Changing Magnetic Field}
\rightarrow
\text{Electric Field}
}
\]

\[
\boxed{
\text{Changing Electric Field}
\rightarrow
\text{Magnetic Field}
}
\]

Therefore:

\[
\boxed{
\text{Electric and magnetic fields can propagate together as waves.}
}
\]

For a plane wave:

\[
\boxed{
\vect{E}
\perp
\vect{H}
\perp
\text{direction of propagation}
}
\]

And the propagation velocity is

\[
\boxed{
v
=
\frac{1}{\sqrt{\mu\epsilon}}
}
\]

For free space,

\[
\boxed{
c
=
\frac{1}{\sqrt{\mu_0\epsilon_0}}
}
\]

Finally, for SDR and signal processing:

\[
\boxed{
\text{RF signal}
\rightarrow
\text{sampling}
\rightarrow
\text{IQ data}
\rightarrow
\text{Fourier/filtering}
\rightarrow
\text{information}
}
\]

That is the bridge between the electromagnetics you're learning and the RF/SDR work you're already doing.

\end{document}
```
